\section{Introduction}
\label{sec:introduction}

Portfolio allocation is a sequential decision problem: holdings selected today affect both next-period wealth and the cost of future rebalancing. The mean--variance paradigm remains a useful benchmark, but estimation error, non-Gaussian marginal behaviour, and changing cross-asset dependence complicate its use in practice \citep{markowitz_portfolio_1952,demiguel_optimal_2009}. These difficulties are particularly consequential when diversification fails in market stress.

Copulas separate marginal dynamics from dependence. A vine copula represents a multivariate copula as a collection of bivariate conditional copulas and can therefore accommodate asymmetric and tail-dependent relationships that are not implied by an elliptical dependence model \citep{bedford_probability_2001,bedford_vinesnew_2002,czado_analyzing_2019}. Reinforcement learning (RL), in turn, provides a natural language for repeated portfolio choice because a policy maps an information set into an action and can account for delayed effects of turnover and wealth accumulation \citep{sutton_reinforcement_2018,lillicrap_continuous_2015}.

Combining the two is appealing but not automatically persuasive. A flexible dependence model does not create return predictability; an RL policy trained on artificial paths can overfit the generator; and a single short backtest cannot establish superiority. The contribution of this paper is consequently methodological and empirical rather than rhetorical. We implement a dependence-aware, leveraged long--short allocation policy and specify the evidence required to assess it credibly.

The implemented framework has three components. First, a training-only Student-$t$ D-vine uses neural forecasts for all 21 unconditional and conditional correlations; monthly snapshots provide a time-indexed dependence summary without rolling refits. Second, a recurrent TD3 agent receives those summaries together with portfolio and monthly market states and produces self-financing long--short weights subject to fixed net and gross exposure. Third, fully synthetic monthly paths are used for pre-training, realised in-sample historical paths are used for fine-tuning, and the final 24 monthly holding periods are excluded from both stages and reserved for evaluation.

This paper makes four deliberately bounded contributions. It (i) documents a concrete dependence-feature state representation; (ii) specifies a training protocol that prevents the held-out period from entering marginal fitting or realised fine-tuning rewards; (iii) distinguishes synthetic distributional fidelity from economic performance; and (iv) lays out a walk-forward, benchmark, ablation, and inference design suitable for testing the framework. It does \emph{not} claim that dependence features, synthetic pre-training, or RL outperform conventional allocation methods before the prescribed experiments are run.

The remainder of the paper reviews the relevant literature, describes the implemented system, documents the data and protocol, provides a results-reporting template and validation plan, and concludes with the changes necessary for a publishable empirical claim.
